\documentclass[10pt,t,english,aspectratio=1610]{beamer}
\usepackage{fontawesome}
\usepackage{graphicx}
\usepackage{array}
\usepackage[normalem]{ulem}
\usepackage{amsfonts,amsmath,amssymb,bm,bbm}
\usepackage{mathrsfs}
% NOTE: sgame is intentionally NOT used -- its `game` environment is
% incompatible with the `array` package (required by colortbl/makecell/multirow)
% and errors with "Missing # inserted in alignment preamble". We use the modern
% TikZ-based `xgames` package instead (loaded after TikZ below).
\usepackage{graphicx,pstricks}
\usepackage{xcolor}
\usepackage{colortbl}
\usepackage{makecell}
\usepackage{tikz,tikzsymbols,gnuplottex}
\usetikzlibrary{decorations.pathreplacing,shapes,positioning,arrows.meta,fit,backgrounds,calc}

% --- xgames: TikZ-based game-theory typesetting (matrixgame, gametree,
% beliefspace, automaton). Not on CTAN; xgames.sty + fikz.sty are vendored at
% the repo root. CC BY-SA 4.0 (xgames) / GPL-3.0 (fikz), (c) Benjamin Bernard.
\usepackage{xgames}

% Custom command for shading rows in a matrixgame
% ---------------------------------------------------------------------------
% \shaderow[<fill options>]{<row>}{<first column>}{<last column>}
% Call this inside the matrixgame, after \payoffmatrix
% Can wrap in a \foreach to shade multiple regions
% Example: \foreach \i in {1,...,5}{ \shaderow[green!50!black]{\i}{\i}{5} }
\newcommand{\shaderow}[4][red!60]{%
  \ifnum#3=1\relax\def\xl{0}\else\pgfmathsetmacro{\xl}{\pmtwi[#3-1]}\fi
  \ifnum#2=1\relax\def\yt{0}\else\pgfmathsetmacro{\yt}{\pmthi[#2-1]}\fi
  \pgfmathsetmacro{\xr}{\pmtwi[#4]}%
  \pgfmathsetmacro{\yb}{\pmthi[#2]}%
  \fill[opacity=0.18, #1] ($(pmt1-tl)+(\xl,-\yt)$) rectangle ($(pmt1-tl)+(\xr,-\yb)$);%
}

\usepackage[english]{babel}
\usepackage[utf8]{inputenc}
\usepackage{appendixnumberbeamer}
\usepackage{datetime2}
\usepackage{booktabs}
\usepackage{setspace}
\usepackage{rotating}
\usepackage{multirow}
\usepackage{listings}

\usepackage{natbib}
\bibliographystyle{unsrtnat}
\renewcommand{\bibfont}{\footnotesize}
\renewcommand{\bibsection}{\vspace{-1em}}

%\usepackage{matlab-prettifier} % For enhanced MATLAB highlighting
\lstset{
    language=matlab,
    % Other options for styling (optional)
    basicstyle=\ttfamily\footnotesize, % Font style
    keywordstyle=\color{blue}, % Keyword color
    commentstyle=\color{green!50!black}, % Comment color
    stringstyle=\color{red!70!black}, % String color
    numbers=left, % Line numbers on the left
    numberstyle=\tiny\color{gray}, % Style for line numbers
    frame=single, % Frame around the listing
    breaklines=true, % Allow line breaking
    captionpos=b, % Caption at the bottom
    tabsize=4 % Tab size
}

% ------------------------------------------------------------------------------
% Use the beautiful metropolis beamer template
% ------------------------------------------------------------------------------
\usepackage[T1]{fontenc}
\usepackage[utf8]{inputenc}
\usepackage{fontawesome}
\usepackage{FiraSans} 
\mode<presentation>
{
  \usetheme[progressbar=foot,background=light]{metropolis} 
  \usecolortheme{default} % or try albatross, beaver, crane, ...
  \usefonttheme{default}  % or try default, serif, structurebold, ...
  \setbeamertemplate{navigation symbols}{}
  \setbeamertemplate{caption}[numbered]
  %\setbeamertemplate{frame footer}{My custom footer}
  \makeatletter
    \setbeamertemplate{frametitle continuation}{}
  \makeatother
} 
\addtobeamertemplate{frametitle continuation}{}{%
  \setbeamertemplate{frametitle continuation}{}%
}
\setbeamertemplate{frametitle continuation}{}

\newenvironment{stepenumerate}{\begin{enumerate}[<+->]}{\end{enumerate}}
\newenvironment{stepitemize}{\begin{itemize}[<+->]}{\end{itemize} }
\newenvironment{stepenumeratewithalert}{\begin{enumerate}[<+-| alert@+>]}{\end{enumerate}}
\newenvironment{stepitemizewithalert}{\begin{itemize}[<+-| alert@+>]}{\end{itemize} }

\newtheorem{question}{Question}
\newtheorem{claim}{Claim}
\newtheorem{proposition}{Proposition}
\newtheorem{remark}{Remark}
\newtheorem{conjecture}{Conjecture}

\definecolor{metrop}{RGB}{29, 44, 44}
\colorlet{GrayLight}{black!15}
\colorlet{GrayMedium}{black!30}
\colorlet{ForestGreen}{green!60!black}

\newenvironment{transitionframe}{
  \setbeamercolor{background canvas}{bg=black!80}
  \begin{frame}}{
    \end{frame}
}

\newcommand{\bsk}{

\bigskip

}

\newcommand{\pd}{\partial}
\newcommand{\RR}{\mathbb{R}}

\newcommand*\hugme[1]{\tikz[baseline=(char.base)]{\node[shape=ellipse,draw,inner sep=0pt] (char) {#1};}}

\newcounter{saveenumi}
\newcommand{\seti}{\setcounter{saveenumi}{\value{enumi}}}
\newcommand{\conti}{\setcounter{enumi}{\value{saveenumi}}}
\resetcounteronoverlays{saveenumi}



\newcommand\dotprod[2]{\langle #1 , #2 \rangle}
\newcommand{\ft}[1]{\widehat #1}
\newcommand{\qabove}[1]{\overset{\text{\large \textbf ?}}{#1}}
\newcommand{\eqae}{\overset{\text{a.e.}}{=}}
\newcommand{\calp}{\mathcal{P}}
\newcommand{\calg}{\mathcal{G}}
\newcommand{\calb}{\mathcal{B}}
\newcommand{\textd}{\text{d}}
\newcommand{\bbr}{\mathbb{R}}
\newcommand{\binm}{\mathbin{M}}
\newcommand{\binc}{\mathbin{C}}
\newcommand{\binb}{\mathbin{B}}
\newcommand{\calc}{\mathcal{C}}
\newcommand{\calh}{\mathcal{H}}
\newcommand{\bfone}{\mathbf{1}}
\newcommand{\bbe}{\mathbb{E}}
\newcommand{\bfle}{\mathbf{e}}
\newcommand{\calf}{\mathcal{F}}
\newcommand{\cala}{\mathcal{A}}
\newcommand{\cale}{\mathcal{E}}
\newcommand{\bbn}{\mathbb{N}}
\newcommand{\cantor}{\calc}
\newcommand{\calY}{\mathcal{Y}}
\newcommand{\textb}{\text{B}}
\newcommand{\calm}{\mathcal{M}}
\newcommand{\bint}{\mathbin{T}}
\newcommand{\ep}{\epsilon}
\newcommand{\bbq}{\mathbb{Q}}
\newcommand{\bbp}{\mathbb{P}}
\newcommand{\cals}{\mathcal{S}}
\newcommand{\emptysequence}{e}
\newcommand{\bbz}{\mathbb{Z}}
\newcommand{\fraka}{\frak{A}}
\newcommand{\frakb}{\frak{B}}
\newcommand{\length}{\text{length}}
\newcommand{\bfn}{\mathbf{N}}
\newcommand{\support}{\text{support}}
\DeclareMathOperator*{\argmax}{arg\,max}
\newcommand{\dom}{\mbox{dom}}
\def\ut{\underline t}
\def\um{\underline m}
\def\PP{\mathbb{P}}
\def\EE{\mathbb{E}}
\def\RR{\mathbb{R}}

% Title page info
\title[Incomplete Information: Static]{Chapter 14: Incomplete Information, Static Games}
\author[ECON 5001]{ECON 5001\\Prof. Paul J. Healy\\Based on: \textit{Game Theory} by Giacomo Bonanno} \color{metrop}
\institute[OSU]{}
\date[]{\vfill {\tiny Updated \today\ at\ \DTMcurrenttime}}

% ---------------------------------------------------------------------
% SCOPE AND NOTATION
%   Builds on deck 11 (states, partitions, beliefs, common prior) and
%   deck 12 (Bayesian updating, sequential rationality). Bonanno's
%   "state-space" approach; types only informally (Ch 16 not covered).
%   Bonanno's Figure 14.3 has the two matrices swapped (see errata).
%   The slides draw it correctly: type a = Figure 14.2, type b =
%   Figure 14.1 (the true game, state beta).
%   Section 14.4 (three players) gets one footnote.
%
% ANSWERS (all checked numerically)
%   One-sided (Bonanno 14.1-14.5). P2 believes type a with prob 2/3.
%     Type a (Fig 14.2): T: (0,3) (3,9)   B: (3,3) (0,0)
%     Type b (Fig 14.1): T: (6,3) (0,9)   B: (3,3) (3,0)   <- true
%     Strategic form (row = P1 "if a, if b"):
%       TT: (2,3) (2,9)   TB: (1,3) (3,6)   BT: (4,3) (0,3)   BB: (3,3) (1,0)
%     BNE: ((T,B), R) and ((B,T), L). Both separating, no pooling.
%     ((T,B),R) at true state beta: play (B,R), not a NE of the true game.
%     Type-by-type check of ((B,T),L): P2 beliefs 2/3 on (a,B), 1/3 on
%       (b,T): L 3, R 3. Type a vs L: B 3 > T 0. Type b vs L: T 6 > B 3.
%   Cournot (Bonanno's numbers, Shell/Speedway): P = 34 - Q, Shell c=6,
%     Speedway c=9 (true, prob 1/3) or c=3 (prob 2/3).
%     Complete info: q_Shell = 31/3 = 10.33, q_Speedway = 22/3 = 7.33.
%     BNE: q_Shell = 9, q_H = 8, q_L = 11.
%   Lemons logic (Akerlof; numbers styled on MobLab's defaults): seller
%     values evenly spread on [500, 5000], buyer values 1.5x. At price p
%     only cars worth <= p are offered; average buyer value
%     .75(500 + p) >= p iff p <= 1500. Only (1500-500)/4500 = 22% of cars
%     can trade, though every trade would be efficient.
%   Two-sided (Bonanno 14.8-14.9). States alpha, beta (game G), gamma
%   (game G'). P1: {alpha}, {beta, gamma} (1/2, 1/2).
%   P2: {alpha, beta} (2/3, 1/3), {gamma}.
%     G:  A: C(1,3) D(0,1)   B: C(0,0) D(3,1)
%     G': A: C(1,0) D(0,2)   B: C(0,0) D(3,1)
%     Common prior (1/2, 1/4, 1/4).
%     BNE: P1 "A at alpha, B at {beta,gamma}", P2 "C at {alpha,beta},
%       D at {gamma}". P2 beliefs left: alpha-A 2/3, beta-B 1/3:
%       C 2, D 1. Right: D 1 > C 0. P1 at alpha: A 1 > B 0.
%       P1 at {beta,gamma}: A 1/2, B 3/2.
%   Worked example (Who Brings the Grill?): benefit 4 each if at least
%   one grill. Cost Low 1, High 5. Independent, Pr(Low) = 1/2 each.
%     4 states LL, LH, HL, HH, prior 1/4 each. True state LL.
%     Pure strategies in the Harsanyi game: 4 each.
%     High never brings (4 - 5 < 0).
%     Unique pure BNE: Low brings, High doesn't (both players).
%       Low: bring 3 vs not 4(1/2) = 2.
%     At LL both bring. True game (both Low): (B,B) (3,3), (B,N) (3,4),
%       (N,B) (4,3), (N,N) (0,0). NE (B,N), (N,B). (B,B) is not.
%     Pr(Low) = .9: Low's not = 3.6 > 3, so the old BNE fails.
%       BNE: Jordan-Low brings and Casey never brings, or vice versa.
% ---------------------------------------------------------------------

% The two possible games in the one-sided example
\newcommand{\gamea}[1][]{%
\begin{matrixgame}[top={L, R}, left={T, B}, player names={1, 2}, #1]
  \payoffmatrix{0,3 & 3,9 \\ 3,3 & 0,0}
\end{matrixgame}}
\newcommand{\gameb}[1][]{%
\begin{matrixgame}[top={L, R}, left={T, B}, player names={1, 2}, #1]
  \payoffmatrix{6,3 & 0,9 \\ 3,3 & 3,0}
\end{matrixgame}}

% Knowledge-belief structure for the one-sided example
\newcommand{\onesided}{%
\begin{tikzpicture}[x=1.4in, y=0.55in]
  \node at (0, 0.75) {$\alpha$ (type $a$)}; \node at (1, 0.75) {$\beta$ (type $b$)};
  \node[anchor=east, text=player1] at (-0.3, 0) {1:};
  \node[anchor=east, text=player2] at (-0.3, -1) {2:};
  \foreach \x in {0,1} {\fill (\x, 0) circle (2pt); \fill (\x, -1) circle (2pt);}
  \draw[player1, thick, rounded corners=7pt] (-0.15, -0.25) rectangle (0.15, 0.25);
  \draw[player1, thick, rounded corners=7pt] (0.85, -0.25) rectangle (1.15, 0.25);
  \draw[player2, thick, rounded corners=7pt] (-0.15, -1.25) rectangle (1.15, -0.75);
  \node[anchor=west, font=\footnotesize] at (0.04, -1) {$\tfrac23$};
  \node[anchor=west, font=\footnotesize] at (1.04, -1) {$\tfrac13$};
  \draw[->, thick] (1, -1.9) -- (1, -1.4); \node[font=\scriptsize] at (1, -2.05) {true state};
\end{tikzpicture}}

% Harsanyi transformation of the one-sided example (Bonanno Fig 14.4)
\newcommand{\harsanyitree}[1][0.82]{%
\begin{gametree}[player names={1, 2}, labels=short, scale=#1]
  \branch[player=0, root=w, label=90] from (0,0) to[v=1.2, h=5.0] {type $a$\ $\frac23$; type $b$\ $\frac13$};
  \branch[player=1, label=180] from (w1) to[v=1.3, h=2.4] {T; B};
  \branch[player=1, label=0] from (w2) to {T; B};
  \branch[player=2, label=none] from (w11) to[v=1.4, h=1.1] {L(0,3); R(3,9)};
  \branch[label=none] from (w12) to {L(3,3); R(0,0)};
  \branch[label=none] from (w21) to {L(6,3); R(0,9)};
  \branch[label=none] from (w22) to {L(3,3); R(3,0)};
  \informset[player=2] (w11) -- (w22);
\end{gametree}}

% Two-sided example (Bonanno Fig 14.9)
\newcommand{\twosidedtree}[1][0.66]{%
\begin{gametree}[player names={1, 2}, labels=short, scale=#1]
  \branch[player=0, root=w, label=90] from (0,0) to[v=1.2, h=5.6] {$\alpha$\ $\frac24$; $\beta$\ $\frac14$; $\gamma$\ $\frac14$};
  \branch[player=1, label=180] from (w1) to[v=1.3, h=2.6] {A; B};
  \branch[player=1, label=none] from (w2) to {A; B};
  \branch[player=1, label=none] from (w3) to {A; B};
  \branch[player=2, label=none] from (w11) to[v=1.4, h=1.25] {C(1,3); D(0,1)};
  \branch[label=none] from (w12) to {C(0,0); D(3,1)};
  \branch[label=none] from (w21) to {C(1,3); D(0,1)};
  \branch[label=none] from (w22) to {C(0,0); D(3,1)};
  \branch[label=none] from (w31) to {C(1,0); D(0,2)};
  \branch[label=none] from (w32) to {C(0,0); D(3,1)};
  \informset[player=1] (w2) -- (w3);
  \informset[player=2] (w11) -- (w22);
  \informset[player=2] (w31) -- (w32);
\end{gametree}}

\begin{document}

\frame{\maketitle}

\begin{frame}{Where We're Going}
So far: the game is \textbf{common knowledge}
\bsk
Now: some player isn't sure \textbf{which game} they're playing
\begin{itemize}
    \item Usually: unsure about someone else's \textbf{payoffs}
    \item ``Is my rival's cost high or low?''
\end{itemize}
\bsk
Tools we have:
\begin{itemize}
    \item Knowledge-belief structures (deck 11): who knows what
    \item Bayesian updating (deck 12): beliefs from strategies
\end{itemize}
\bsk
New tool: the \textbf{Harsanyi transformation}
\vspace{-0.3in}
\end{frame}

% ---- MobLab run: Market for Lemons ------------------------------------
% Markets topic, two-sided version. Fixed pairs, Fixed Role. Seller
% values 500 to 5000 (default); Beta = 150 (buyer value = 1.5 x seller
% value). MobLab doesn't show Beta to students: the slide announces it.
% Asymmetric information treatment. Don't announce the number of
% periods. Results come back in "Back to MobLab".
% ---------------------------------------------------------------------
\begin{frame}{MobLab: ``Market for Lemons''}
\begin{itemize}
    \item MobLab time!
    \item You're a \textbf{seller} or a \textbf{buyer} of a used car
    \item Seller knows what the car is worth to them: between \$500 and \$5{,}000
    \item Seller names a price. Buyer accepts or rejects
    \item Buyer doesn't see the car's value until after
    \item The car is worth 50\% more to the buyer than to the seller
\end{itemize}
\vfill
\centering
Feel free to talk! And play it the way you'd actually play it.
\vspace{-0.2in}
\end{frame}

\begin{frame}{Imperfect vs.\ Incomplete Information}
\textbf{Imperfect information} (Chapter 4)
\begin{itemize}
    \item You don't know what others \textbf{did} earlier
    \item Information sets with more than one node
\end{itemize}
\bsk
\textbf{Incomplete information} (now)
\begin{itemize}
    \item You don't know what \textbf{game} you're playing
    \item Someone's payoffs (or actions, or...) are uncertain
\end{itemize}
\bsk
\textbf{Ex:} someone says you owe them money and threatens to sue
\begin{itemize}
    \item Will they really sue? Depends on their \textbf{type}: aggressive or not
    \item You don't know their type. They do
\end{itemize}
\vspace{-0.3in}
\end{frame}

% ==============================================================
% 14.2 ONE-SIDED INCOMPLETE INFORMATION
% ==============================================================

\begin{frame}[standout]
One-Sided Incomplete Information
\end{frame}

\begin{frame}{Example: Two Possible Games}
\begin{columns}[T]
\begin{column}{0.48\textwidth}
\centering
\textbf{P1 is type $a$}\\[4pt]
\gamea
\end{column}
\begin{column}{0.48\textwidth}
\centering
\textbf{P1 is type $b$} (the true game)\\[4pt]
\gameb
\end{column}
\end{columns}
\bsk
\begin{itemize}
    \item Only P1's payoffs differ
    \item P1 knows which game it is
    \item P2 thinks: type $a$ with prob.\ $\tfrac23$, type $b$ with prob.\ $\tfrac13$
    \item All of this is common knowledge
\end{itemize}
\vfill
{\footnotesize Bonanno Figures 14.1--14.2.}
\vspace{-0.2in}
\end{frame}

\begin{frame}{The Knowledge-Belief Structure}
\centering
\onesided
\bsk
\raggedright
Each state comes with a \textbf{game}: $\alpha\to$ type $a$'s game, $\beta\to$ type $b$'s game
\bsk
What's common knowledge? (Constant across states!)
\begin{itemize}
    \item \underline{\hspace{1.5in}} \quad \underline{\hspace{1.5in}} \quad \underline{\hspace{1.5in}}
\end{itemize}
\vfill
{\footnotesize Bonanno Figure 14.3 draws the two games swapped. This is the correct version.}
\vspace{-0.2in}
\end{frame}

\begin{frame}{The Harsanyi Transformation}
\begin{columns}[c]
\begin{column}{0.55\textwidth}
\centering
\harsanyitree
\end{column}
\begin{column}{0.43\textwidth}
\begin{enumerate}
    \item \textbf{Nature} picks the state
    \begin{itemize}
        \item Prob.\ $=$ P2's beliefs
    \end{itemize}
    \item P1 learns it, chooses T or B
    \item P2 chooses L or R \textbf{without} seeing either
\end{enumerate}
\end{column}
\end{columns}
\bsk
Now it's a regular game! Imperfect information, but \textbf{complete}
\vspace{-0.3in}
\end{frame}

\begin{frame}{John Harsanyi (1920--2000)}
\begin{columns}
\begin{column}{0.40\textwidth}
    \begin{center}
        \IfFileExists{LectureSlides/graphics/history/HarsanyiMug.png}{%
        \includegraphics[height=1.3in]{LectureSlides/graphics/history/HarsanyiMug.png}}{%
        \fbox{\parbox[c][1.3in][c]{1in}{\centering\footnotesize (photo)}}}
    \end{center}
\end{column}
\begin{column}{0.60\textwidth}
\begin{itemize}
    \item Budapest. Same high school as von Neumann
    \item 1944: escaped a deportation train at the station
    \item 1950: snuck across a marsh into Austria
    \item Sydney: factory job by day, economics at night
    \item Games of incomplete information: 1967--68, three papers
    \item 1994 Nobel Prize, with Nash and Selten
\end{itemize}
\end{column}
\end{columns}
\vfill
{\footnotesize Not in the book. Harsanyi's Nobel autobiography (nobelprize.org).}
\vspace{-0.2in}
\end{frame}

\begin{frame}{Strategies and the Strategic Form}
P1 has two information sets (one per type) $\Rightarrow$ strategy $=$ (choice if $a$, choice if $b$)
\bsk
\centering
\only<1>{\begin{matrixgame}[top={L, R}, left={TT, TB, BT, BB}, player names={Player 1, Player 2}, w=1.9]
  \payoffmatrix{2,3 & 2,9 \\ {}, & {}, \\ {}, & {}, \\ 3,3 & 1,0}
\end{matrixgame}}%
\only<2->{\begin{matrixgame}[top={L, R}, left={TT, TB, BT, BB}, player names={Player 1, Player 2}, w=1.9, br=<3->]
  \payoffmatrix{2,3 & 2,9 \\ 1,3 & 3,6 \\ 4,3 & 0,3 \\ 3,3 & 1,0}
\end{matrixgame}}
\bsk
\raggedright
Payoffs are \textbf{expected} payoffs. \ \textbf{Ex:} TT vs.\ L, P1 gets $\tfrac23(0)+\tfrac13(6)=2$\\[4pt]
\onslide<1>{Fill in TB and BT}
\vspace{-0.3in}
\end{frame}

\begin{frame}{Bayesian Nash Equilibrium}
\begin{block}{Definition}
A \textbf{Bayesian Nash equilibrium} (BNE) is a NE of the game from the Harsanyi transformation
\end{block}
\bsk
``Bayesian'' just reminds us it came from incomplete information
\bsk
Our example: pure BNE are \underline{\hspace{1.1in}} and \underline{\hspace{1.1in}}
\bsk
\onslide<2->{\textbf{Fact:} the Harsanyi game is finite $\Rightarrow$ a BNE always exists (maybe mixed)}
\vspace{-0.3in}
\end{frame}

\begin{frame}{What Does a BNE Mean?}
Take the BNE $((T,B),R)$. The true state is $\beta$ (type $b$)
\begin{itemize}
    \item What actually gets played? \underline{\hspace{0.8in}}
    \item Is that a NE of the true game? \underline{\hspace{0.8in}}
\end{itemize}
\bsk
\textbf{Lesson:} a BNE need not be a NE of the game actually being played
\bsk
\onslide<2->{But P1 \textbf{knows} she's type $b$. Why plan for type $a$??
\begin{itemize}
    \item Her strategy $(T,B)$ $=$ what she actually does ($B$)\\
    $+$ what P2 \textbf{thinks} type $a$ would do ($T$)
\end{itemize}}
\vspace{-0.3in}
\end{frame}

\begin{frame}{Pooling vs.\ Separating}
\begin{block}{Definition (Bonanno 14.2.1)}
In a pure BNE, the informed player is
\begin{itemize}
    \item \textbf{Pooling} if every type makes the \textbf{same} choice
    \item \textbf{Separating} if different types make \textbf{different} choices
\end{itemize}
\end{block}
\bsk
Our BNE: $((T,B),R)$ is \underline{\hspace{1in}} \quad $((B,T),L)$ is \underline{\hspace{1in}}
\bsk
\onslide<2->{Separating: P1's choice would \textbf{reveal} her type, if P2 could see it}
\vspace{-0.3in}
\end{frame}

\begin{frame}{Shortcut: Check Type by Type}
\begin{columns}[c]
\begin{column}{0.52\textwidth}
\centering
\harsanyitree[0.72]
\end{column}
\begin{column}{0.46\textwidth}
Is $((B,T),L)$ a BNE?
\bsk
\textbf{P2's beliefs} (Bayes):\\
(type $a$, B): \underline{\hspace{0.4in}} \ (type $b$, T): \underline{\hspace{0.4in}}
\bsk
P2: L gives \underline{\hspace{0.3in}} \ R gives \underline{\hspace{0.3in}}
\bsk
\textbf{Type $a$} vs.\ L: \ T \underline{\hspace{0.25in}} \ B \underline{\hspace{0.25in}}\\[4pt]
\textbf{Type $b$} vs.\ L: \ T \underline{\hspace{0.25in}} \ B \underline{\hspace{0.25in}}
\end{column}
\end{columns}
\bsk
This is WSE logic: Bayesian updating $+$ best replies. Bonanno uses it a lot
\vfill
{\footnotesize WSE $\Rightarrow$ NE (Theorem 11.3.1), so a WSE of the Harsanyi game gives a BNE.}
\vspace{-0.2in}
\end{frame}

% ==============================================================
% COURNOT WITH UNKNOWN COSTS
% ==============================================================

\begin{frame}{Shell vs.\ Speedway, Again}
Cournot (deck 04): Shell and Speedway choose gallons $q_1$, $q_2$. \ $P=34-Q$
\begin{itemize}
    \item Shell's cost: \$6/gallon. Common knowledge
    \item Speedway's cost: really \$9/gallon (High)
    \item But Shell thinks: High ($\tfrac13$) or a new low-cost tech, \$3 (Low, $\tfrac23$)
    \item Speedway knows its own cost
\end{itemize}
\bsk
\textbf{Benchmark}, complete info (both know $c_2=9$):
\begin{itemize}
    \item Shell: $\max_{q_1}\ (34-q_1-q_2)q_1-6q_1$ \quad Speedway: $\max_{q_2}\ (34-q_1-q_2)q_2-9q_2$
    \item $q_1=$ \underline{\hspace{0.5in}} \qquad $q_2=$ \underline{\hspace{0.5in}}
\end{itemize}
\vfill
{\footnotesize Bonanno \S14.2. Uses calculus, like deck 04.}
\vspace{-0.2in}
\end{frame}

\begin{frame}{Cournot BNE}
Strategy for Speedway: one output \textbf{per type}: $q_H$ and $q_L$
\bsk
\textbf{High Speedway:} $\max_{q_H}\ (34-q_1-q_H)q_H-9q_H \ \Rightarrow$ \underline{\hspace{1.5in}}
\bsk
\textbf{Low Speedway:} $\max_{q_L}\ (34-q_1-q_L)q_L-3q_L \ \Rightarrow$ \underline{\hspace{1.5in}}
\bsk
\textbf{Shell}, who doesn't know the type:
$$\max_{q_1}\ \tfrac13\big[(34-q_1-q_H)q_1-6q_1\big]+\tfrac23\big[(34-q_1-q_L)q_1-6q_1\big]$$
$\Rightarrow$ \underline{\hspace{1.5in}}
\bsk
\onslide<2->{Solve: \ $q_1=$ \underline{\hspace{0.4in}} \quad $q_H=$ \underline{\hspace{0.4in}} \quad $q_L=$ \underline{\hspace{0.4in}}}
\vspace{-0.3in}
\end{frame}

\begin{frame}{Compare}
\centering
\renewcommand{\arraystretch}{1.6}
\begin{tabular}{|l|c|c|}
\hline
 & Complete info & Incomplete info \\
\hline
Shell & 10.33 & \hspace{0.8in} \\
Speedway (really High) & 7.33 & \\
\hline
\end{tabular}
\bsk
\raggedright
\onslide<2->{Shell fears a low-cost (aggressive) rival $\Rightarrow$ Shell produces \underline{\hspace{0.6in}}\\[4pt]
So High Speedway faces less competition $\Rightarrow$ produces \underline{\hspace{0.6in}}}
\bsk
\onslide<3->{\textbf{Lesson:} beliefs about your type change what others do, even when they're wrong}
\vspace{-0.3in}
\end{frame}

\begin{frame}{This Is Everywhere}
Bayesian games: you know your own type, not theirs
\begin{itemize}
    \item \textbf{Auctions:} you know your value for the item, not others'
    \item \textbf{Used cars:} the seller knows if it's a lemon
    \item \textbf{Insurance:} you know how risky you are. The insurer doesn't
    \item \textbf{Job interviews:} you know how good you are
    \item \textbf{Poker:} you know your cards
    \item \textbf{Negotiation:} you know your walk-away price
\end{itemize}
\bsk
Real versions: many types, continuous types, many players, repeated play
\begin{itemize}
    \item Same idea: each type best-responds to beliefs about others' types
\end{itemize}
\vspace{-0.3in}
\end{frame}

\begin{frame}{Back to MobLab: Lemons}
\centering
\renewcommand{\arraystretch}{1.8}
\begin{tabular}{|l|c|}
\hline
\% of cars sold & \hspace{1in} \\
\hline
\% of buyers who lost money & \\
\hline
\end{tabular}
\bsk
\raggedright
\textbf{Akerlof's logic:} say seller values are spread evenly from \$500 to \$5{,}000
\begin{itemize}
    \item At price $p$: only cars worth $\le p$ to their seller get offered
    \item Avg.\ value to the buyer: $1.5\times\frac{500+p}{2}$. Worth paying $p$ only if $p\le$ \underline{\hspace{0.5in}}
    \item Only \underline{\hspace{0.4in}}\% of cars trade, though \textbf{every} trade would help both!
\end{itemize}
\vfill
{\footnotesize Not in the book. Akerlof (1970, \textit{QJE}). 2001 Nobel Prize.}
\vspace{-0.2in}
\end{frame}

% ==============================================================
% 14.3 TWO-SIDED INCOMPLETE INFORMATION
% ==============================================================

\begin{frame}[standout]
Two-Sided Incomplete Information
\end{frame}

\begin{frame}{Example: Both Players Unsure}
\begin{columns}[c]
\begin{column}{0.58\textwidth}
\centering
\begin{tikzpicture}[x=0.95in, y=0.55in]
  \foreach \x/\s/\g in {0/\alpha/G, 1/\beta/G, 2/\gamma/G'} {\node at (\x, 0.75) {$\s$}; \node[font=\footnotesize] at (\x, 1.2) {game $\g$};}
  \node[anchor=east, text=player1] at (-0.3, 0) {1:};
  \node[anchor=east, text=player2] at (-0.3, -1) {2:};
  \foreach \x in {0,1,2} {\fill (\x, 0) circle (2pt); \fill (\x, -1) circle (2pt);}
  \draw[player1, thick, rounded corners=7pt] (-0.18, -0.25) rectangle (0.18, 0.25);
  \draw[player1, thick, rounded corners=7pt] (0.82, -0.25) rectangle (2.18, 0.25);
  \draw[player2, thick, rounded corners=7pt] (-0.18, -1.25) rectangle (1.18, -0.75);
  \draw[player2, thick, rounded corners=7pt] (1.82, -1.25) rectangle (2.18, -0.75);
  \foreach \x/\y/\b in {1/0/\tfrac12, 2/0/\tfrac12, 0/-1/\tfrac23, 1/-1/\tfrac13}
    {\node[anchor=west, font=\footnotesize] at (\x+0.04, \y) {$\b$};}
  \draw[->, thick] (0, -1.9) -- (0, -1.4); \node[font=\scriptsize] at (0, -2.05) {true state};
\end{tikzpicture}
\end{column}
\begin{column}{0.4\textwidth}
\centering
Game $G$\\[2pt]
\begin{matrixgame}[top={C, D}, left={A, B}, player names={1, 2}, w=1.3, h=0.65]
  \payoffmatrix{1,3 & 0,1 \\ 0,0 & 3,1}
\end{matrixgame}\\[6pt]
Game $G'$\\[2pt]
\begin{matrixgame}[top={C, D}, left={A, B}, player names={1, 2}, w=1.3, h=0.65]
  \payoffmatrix{1,0 & 0,2 \\ 0,0 & 3,1}
\end{matrixgame}
\end{column}
\end{columns}
\bsk
At $\alpha$: both know it's $G$. But P2 isn't sure P1 knows!
\vspace{-0.3in}
\end{frame}

\begin{frame}{Which Probabilities for Nature?}
One-sided: use the uninformed player's beliefs. Two-sided: \textbf{two} sets of beliefs!
\begin{itemize}
    \item P1 puts $\tfrac12$ on $\beta$. P2 puts $\tfrac13$ on $\beta$. Which one?
\end{itemize}
\bsk
\onslide<2->{Neither! Use the \textbf{common prior} (deck 11): one $p$ that gives both players' beliefs
\vspace{0.6in}

$p(\alpha)=$ \underline{\hspace{0.4in}} \quad $p(\beta)=$ \underline{\hspace{0.4in}} \quad $p(\gamma)=$ \underline{\hspace{0.4in}}}
\bsk
\onslide<3->{\textbf{No common prior?} Then no Harsanyi transformation (Bonanno Exercise 14.6)}
\vspace{-0.3in}
\end{frame}

\begin{frame}{The Harsanyi Game}
\centering
\twosidedtree
\bsk
\raggedright
P1 knows if the state is $\alpha$. P2 knows if the state is $\gamma$
\vfill
{\footnotesize Bonanno Figure 14.9.}
\vspace{-0.2in}
\end{frame}

\begin{frame}{Checking a BNE}
\centering
\twosidedtree[0.5]
\bsk
\raggedright
\begin{columns}[T]
\begin{column}{0.48\textwidth}
Candidate:\\
P1: A at $\alpha$, B at $\{\beta,\gamma\}$\\
P2: C at $\{\alpha,\beta\}$, D at $\{\gamma\}$
\end{column}
\begin{column}{0.5\textwidth}
P2's left info set. Beliefs: \underline{\hspace{0.8in}}\\[3pt]
P2: C gives \underline{\hspace{0.3in}}, D gives \underline{\hspace{0.3in}}\\[3pt]
P1 at $\{\beta,\gamma\}$: A \underline{\hspace{0.3in}}, B \underline{\hspace{0.3in}}
\end{column}
\end{columns}
\vfill
{\footnotesize Three or more players (\S14.4): same steps, more bookkeeping.}
\vspace{-0.2in}
\end{frame}

% ==============================================================
% WORKED EXAMPLE
% ==============================================================
% Answers are in the header comment.
% ---------------------------------------------------------------------

\begin{frame}[standout]
Let's Work One Out
\end{frame}

\begin{frame}{Who Brings the Grill?}
Jordan and Casey tailgate before the game. Each chooses: \textbf{Bring} a grill or \textbf{Not}
\begin{itemize}
    \item If at least one grill shows up: each gets 4. Otherwise 0
    \item Bringing costs 1 if you own a grill (\textbf{Low} type), 5 if you'd rent one (\textbf{High} type)
    \item You know your own type. Each is Low with prob.\ $\tfrac12$, independently
    \item Actually both are Low
\end{itemize}
\bsk
\textbf{Ex:} Jordan Low brings, Casey doesn't: Jordan gets $4-1=3$, Casey gets 4
\vfill
{\footnotesize Choices are simultaneous.}
\vspace{-0.2in}
\end{frame}

\begin{frame}{Who Brings the Grill?}
\begin{enumerate}
    \item Draw the knowledge-belief structure (4 states). What's the common prior? \vspace{0.3in}
    \item How many pure strategies does each player have in the Harsanyi game? \vspace{0.2in}
    \item Show that a High type never brings a grill \vspace{0.2in}
    \item Find all pure BNE. Pooling or separating? \vspace{0.3in}
    \item What happens at the true state? Is it a NE of the true game? \vspace{0.3in}
    \item Now Pr(Low) $=0.9$. Is your BNE still a BNE? If not, find one \vspace{0.1in}
\end{enumerate}
\vspace{-0.3in}
\end{frame}

\begin{frame}{What Have We Learned, Charlie Brown?}
\begin{itemize}
    \item \textbf{Incomplete info}: unsure which game. \textbf{Imperfect}: unsure what others did
    \item Model it with a knowledge-belief structure: one game per state
    \item \textbf{Harsanyi transformation}: Nature picks the state using the common prior
    \item \textbf{BNE} $=$ NE of the Harsanyi game
    \begin{itemize}
        \item Strategy: one choice per type. Check type by type
        \item Pooling vs.\ separating
    \end{itemize}
    \item A BNE may not be a NE of the true game
    \item Two-sided: needs a common prior
\end{itemize}
\vfill
\vspace{-0.4in}
\end{frame}

\begin{frame}{Reading \& Practice}
\begin{itemize}
    \item Bonanno, Chapter 14
    \begin{itemize}
        \item \S14.1--14.3: read. \S14.4: skim
        \item Figure 14.3 has the two games swapped (see errata)
    \end{itemize}
    \item Exercises: \S14.5.1 (14.1, 14.2), \S14.5.2 (14.3--14.5), \S14.5.3 (14.6)
    \item Skip 14.7 (three players) and 14.8 (challenging)
    \item Solutions in \S14.6. Try them first
\end{itemize}
\vfill
\end{frame}

\end{document}
